%% ============================================================
%% en/sections/08_passione.tex
%% Optional layer: Pasión de las Pasiones
%% ============================================================

\chapter{The Passion of the Vespers}
\markboth{The Passion of the Vespers}{The Passion of the Vespers}

\lettrine[lines=3]{T}{he} history of the Sicilian Vespers is, among other
things, a history of passions. Impossible loves between enemy families.
Betrayals born of love and ending in blood. Parents who sacrifice their
children to the cause. Rivals who hate each other with the same intensity
they once loved.

This chapter presents an \textbf{optional dramatic layer} inspired by
Brandon León-Gambetta's \textit{Pasión de las Pasiones}. It does not
replace the FitD system — it runs alongside it. Each character may have,
in addition to their playbook, a \textbf{Passion Role}: a second identity
layer that defines the nature of their personal drama and guarantees that,
sooner or later, that drama will explode.

\nota[GM Note — When to Use This Layer]{
  The Passion layer is optional and activates table by table. It works
  best with players who enjoy personal drama alongside political action.
  If your table is more strategy-oriented than relationship-oriented,
  you can ignore it entirely without losing anything essential. If you
  want your players to cry — and laugh — activate it.
}

\secsep

\section{How It Works}

\subsection{The Passion Role}

At the start of the session, each player who wants to participate in
the layer chooses a Passion Role from the list in this chapter. The Role
is independent of the playbook: a Chatelaine can be The Impossible Love;
a Preaching Friar can be The Betrayer.

Each Role has:
\begin{itemize}[nosep]
  \item \textbf{The Torment} --- the emotional tension that defines the character
  \item \textbf{The Trigger} --- the situation that makes the drama explode
  \item \textbf{The Promise} --- something that \textit{must} happen before
    the end of the session (the GM guarantees it narratively)
  \item \textbf{The Moment} --- a special scene the player can request once only
  \item \textbf{The Passion Score} --- a resource separate from Stress
\end{itemize}

\subsection{The Passion Score}

Each character with a Passion Role has \textbf{5 Passion boxes}.
Passion is distinct from Stress: it represents the emotional intensity
of the moment, not physical or psychological damage.

\regola[Gaining Passion]{
  You gain 1 Passion when:
  \begin{itemize}[nosep]
    \item your Trigger activates
    \item someone you love is in danger because of you
    \item you must choose between the cause and your heart
    \item someone discovers one of your secrets
  \end{itemize}
}

\regola[Spending Passion]{
  You may spend Passion (1 box) to:
  \begin{itemize}[nosep]
    \item add +1d to any roll involving emotions, relationships,
      or protecting someone you love
    \item declare a \textbf{Dramatic Moment}: your character says or does
      something memorable --- the GM must find narrative space for it
    \item refuse the consequences of a roll relating to your Torment,
      replacing them with an emotional consequence instead of a physical one
  \end{itemize}
}

\regola[Passion at Maximum]{
  When all 5 boxes are full, the character enters \textbf{Passion Crisis}:
  they must act according to their Torment, even if counterproductive.
  After the Crisis, all boxes clear. A well-played Passion Crisis
  automatically adds +1 to the faction score --- because personal drama,
  in this game, also moves history.
}

\secsep

\section{The Passion Roles}

\subsection{The Protagonist}
\textit{Someone decided your destiny. You decided otherwise.}

\textbf{Torment:} You have been defined by what others wanted from you
--- your name, your family, your God, your cause. You are discovering,
right now, in the middle of this crisis, that perhaps you are something
different. The problem is you do not yet know what.

\textbf{Trigger:} Someone tells you who you are or what you must do.
Every time this happens, gain 1 Passion.

\textbf{Promise:} Before the session ends, you will make a choice that
contradicts what everyone expects from you. It will not be the right
choice. It will be yours.

\textbf{Moment:} Once only, you may declare: \textit{``Enough.''} Your
character stops, turns, and says or does something no one expected.
Roll 0 dice --- the result is automatically a full success, but the
narrative consequence is irreversible.

\textbf{Adapt to 1282:} The Traitor's Wife who decides to stop hiding.
The Son who disobeys his father at the crucial moment. The Abbess who
opens the convent doors knowing what will follow.

\filetto

\subsection{The Impossible Love}
\textit{You love them. You cannot love them. You keep loving them.}

\textbf{Torment:} You love someone you should not --- for class, faction,
faith, or the blood spilled between your families. It is not a choice.
It is not a weakness. It is simply true, and the world leaves no room
for this truth.

\textbf{Trigger:} Every time that person is in the same scene as you,
gain 1 Passion. Every time you put them in danger with your actions,
gain 2 Passion.

\textbf{Promise:} Before the session ends, you will have a moment alone
--- even if brief, even if interrupted. The GM guarantees it. What
happens in that moment is up to you.

\textbf{Moment:} Once only, you may cross any line --- faction, danger,
common sense --- to protect or reach that person. Roll normally, but
whatever the result, you get there. You pay the cost after.

\textbf{Adapt to 1282:} An Angevin knight who loves a Sicilian woman.
A friar in love with someone their vows forbid. A Conspirator in love
with someone in the wrong faction.

\nota[Cross-Table Note]{
  The Impossible Love works best when the beloved is a PC at another
  table. The coordinating GM can manage shared moments as special events
  between Conjunction Moments. If this is not logistically possible,
  the beloved may be a significant NPC.
}

\filetto

\subsection{The Betrayer}
\textit{You have a secret. The secret will destroy everything. You cannot stop carrying it.}

\textbf{Torment:} You are doing something --- or have done something ---
that would betray everything others expect from you. You are not
necessarily a villain: perhaps you do it for love, for fear, for a debt
impossible to repay. But if it comes out, nothing will be the same.

\textbf{Trigger:} Every time someone gets close to the truth, gain 1
Passion. Every time you actively lie about what you are doing, gain
1 Passion.

\textbf{Promise:} The secret will come out before the session ends.
Not necessarily in front of everyone --- but at least in front of one
person who should not have known. The GM guarantees it. How it happens
depends on your actions.

\textbf{Moment:} Once only, you may choose to voluntarily confess your
secret to someone. If you do, clear all Passion boxes and gain +2 dice
on your next roll --- whatever it is. The confession does not change
the facts. It changes who you are.

\textbf{Adapt to 1282:} The Spy playing a double game even with the
Conspirators. The vassal informing the French to protect their family.
The Canon who knows something devastating and has never said it.

\filetto

\subsection{The Parental Figure}
\textit{You gave everything to protect someone. Tonight it may not be enough.}

\textbf{Torment:} Someone depends on you --- a child, a nephew, an entire
community, a young person you raised as your own. You built your whole
life around this protection. And now the situation asks you to choose
between them and the cause.

\textbf{Trigger:} Every time the person you protect is in danger, gain
2 Passion. Every time you must choose between them and the cause, gain
1 Passion.

\textbf{Promise:} Before the session ends, the person you protect will
be in real danger --- not because of you, but because of events. How
you save them (or do not) is the heart of your story.

\textbf{Moment:} Once only, you may place yourself between your beloved
and the danger. You automatically absorb all the harm they would have
suffered. No roll. Then roll Resist to see what it cost you.

\textbf{Adapt to 1282:} The Fishwife and the Boy. The Old One and
whoever reminds them of someone. The Bishop and their flock.
The Baron and the Son who wants to fight.

\filetto

\subsection{The Rival}
\textit{They want what you have. Or perhaps it is the reverse. At this point neither of you remembers.}

\textbf{Torment:} There is someone --- another PC or a significant NPC
--- with whom you have a history. Perhaps you were friends, lovers,
siblings. Now you are adversaries. Not pure hatred: something more
complicated, where respect and resentment mix in proportions that change
every scene.

\textbf{Trigger:} Every time your Rival gets something you wanted,
gain 1 Passion. Every time you surpass them, gain 1 Passion (yes,
victory costs as much as defeat).

\textbf{Promise:} Before the session ends, you will have a direct
confrontation --- not necessarily violent. Something will be said
that could not remain unsaid.

\textbf{Moment:} Once only, you may declare you are helping your Rival
instead of obstructing them. This does not change your relationship ---
it complicates it. Gain 3 Passion and the Rival owes you something
they cannot refuse.

\textbf{Adapt to 1282:} Two barons with an old score to settle.
The Conspirator and the Spy who once worked together. The Preaching
Friar and a preacher from the opposing side.

\filetto

\subsection{The Innocent Bystander}
\textit{You did not want to be here. Now you cannot leave.}

\textbf{Torment:} You ended up in this story by chance, by misfortune,
or because someone you loved dragged you in. You have no grand ideals.
You have no plan. You only have the growing feeling that tonight will
change everything --- and that you will not be the same person after.

\textbf{Trigger:} Every time you see something you should not have seen,
gain 1 Passion. Every time someone you thought was good does something
terrible, gain 1 Passion.

\textbf{Promise:} Before the session ends, you will do something brave.
Not because you are a hero --- but because you could not do otherwise.
The GM guarantees the moment will arise.

\textbf{Moment:} Once only, you may declare you were in the right place
at the right time. Something that was about to go very wrong does not ---
because you were there. No roll. The GM integrates it narratively.

\textbf{Adapt to 1282:} The Boy who does not yet understand what is
happening. The Pilgrim who should have left yesterday. The Arab Physician
who only wanted to treat someone.

\filetto

\subsection{The Martyr}
\textit{You already know how this ends. You carry on anyway.}

\textbf{Torment:} You have seen enough to understand that your path leads
toward something painful --- perhaps death, perhaps the loss of everything
you love, perhaps simply the end of who you were. You are not trying to
avoid it. You are trying to make it matter.

\textbf{Trigger:} Every time someone tries to stop or protect you, gain
1 Passion. Every time you choose the harder path when an easier one
exists, gain 1 Passion.

\textbf{Promise:} Before the session ends, you will be offered a way out.
You may choose whether to take it or not.

\textbf{Moment:} Once only, you may do something irreversible that
guarantees the cause succeeds --- at a personal cost you define yourself.
The GM cannot deny the effect. The cost cannot be denied by anyone,
including you.

\textbf{Adapt to 1282:} The Noble in Exile who knows they will return
home only to die. The Preaching Friar who preaches knowing they will be
arrested. The Widow who has nothing left to lose.

\filetto

\subsection{The Redeemed}
\textit{You have done things you are ashamed of. Tonight you might do something different.}

\textbf{Torment:} Your past weighs on you. Perhaps you collaborated with
the Angevins, perhaps you betrayed someone, perhaps you simply looked
away too many times. You are not seeking others' forgiveness --- you are
trying to find out whether you are still capable of deserving it.

\textbf{Trigger:} Every time someone treats you as though you are still
the person you were, gain 1 Passion. Every time you have the chance to
do the right thing and actually do it, gain 1 Passion.

\textbf{Promise:} Before the session ends, your past will return ---
someone, something, a consequence of what you did. How you face it
defines who you are now.

\textbf{Moment:} Once only, you may declare you are doing something your
old self would never have done. Roll with advantage (+1d) and the GM
must acknowledge the difference narratively.

\textbf{Adapt to 1282:} The Deserter. The Traitor's Wife. The Merchant
who sold information once and does not want to again.

\secsep

\section{Passion Connections}

At the start of the session, after choosing Roles, each player declares
a \textbf{Passion Connection}: a specific relationship with another PC
(at the same table, or at another table if the coordinating GM permits)
that will fuel the drama.

\begin{itemize}[nosep]
  \item The Connection may be love, hatred, debt, blood, or shared secret
  \item Every time the Connection is relevant in a scene, both characters
    involved gain 1 Passion
  \item If the Connection is broken during the session (betrayal, death,
    definitive separation), both gain 3 Passion simultaneously
\end{itemize}

\nota[Suggestion --- Pre-built Connections]{
  To speed up setup, you may assign pre-built Connections based on
  playbook Opening Bonds. The Knight and the Son are already in a
  mentor/student relationship perfect for Parental Figure + Protagonist.
  The Conspirator and the Spy are already a natural Rival pair.
}

\nota[Optional Prop --- The Connection Ribbon]{
  In live play, each Connection can be represented by a ribbon or cord
  that physically ties the two players together (see ``Optional Physical
  Components'' in the rules chapter). The colour may encode the nature
  of the bond --- love, hatred, debt, blood, secret --- but a neutral
  ribbon works too. When the Connection breaks, the ribbon is cut: the
  gesture is the rule.
}

\subsection{To Bind and to Loose --- the Clergy's Action}

The Church claims the power of \textit{ligare et solvere}: to bind and
to loose, in heaven and on earth. At the Clergy table, this is not only
theology.

\regola[To Bind and to Loose]{
  A character of the \textbf{Clergy} faction with a Passion Role may,
  \textbf{once per act}, spend Passion to act upon \textit{other people's}
  Connections --- never their own:
  \begin{itemize}[nosep]
    \item \textbf{Bind} (spend 2 Passion) --- declare a new Connection
      between two other characters (PCs or significant NPCs). The two
      bound cannot refuse the bond, but they choose its nature. A ribbon
      is stretched between them.
    \item \textbf{Loose} (spend 3 Passion) --- break an existing
      Connection between two other characters. This counts as a broken
      Connection: both bound characters gain 3 Passion. The ribbon is cut.
  \end{itemize}
  In both cases the Clergy must justify the act in the fiction: a
  blessing, a confession heard, a marriage, an excommunication, a
  revealed truth. Without the scene, there is no power.
}

\nota[GM Note --- Balance]{
  To Bind and to Loose is powerful: it gives the Clergy leverage over the
  drama of every other table. This is intentional --- the Church of 1282
  \textit{was} this. But if a Clergy player abuses it, remind them that
  every imposed bond breeds resentment: characters bound or loosed against
  their will have every reason to make the Clergy pay for it, in scene
  and in the faction score.
}

\secsep

\section{The Passion Scene}

Once per session --- typically around the second Conjunction Moment
(The Crisis) --- the GM may declare a \textbf{Passion Scene}: a brief
scene outside the main action flow in which two characters (with a
Passion Connection between them) find themselves alone for 5--10 minutes
of play.

No mechanics. No rolls. Just the scene.

At the end of the Passion Scene, both players declare what has changed
in their character. They may gain or spend Passion freely. The GM takes
note and integrates the consequences into the finale.

\nota[Design Note --- Why This Works]{
  The Passion layer does not slow down the political game --- it grounds
  it. When players have emotional investment in personal relationships,
  political choices gain weight. The Baron who must choose between the
  cause and their child is more interesting than the Baron choosing
  between two strategic options. Personal drama \textit{is} strategy,
  in Vespri 1282 --- because it was in the real history too.
}

\filettodoppio
